\section{A global scalar inverse with all low cap modes}
\label{sec:global-inverse}

Fix $k\ge2$, $0<\alpha<1$, and
$0<\eta<\min(\mu,1)$, where $\mu$ is the rate in
Section~\ref{sec:initial-metrics}. Put $\delta=2+\eta$.
Choose a small fixed radius $r_0$ with disjoint node
neighbourhoods of radius $6r_0$, and take $R<r_0/10$.
For fixed $R$ all cap cores lie inside $r_i<R$ once
$p$ is sufficiently small. Write
\[
 P_p^R=\Lambda_{\widetilde\omega_p^R}
                           \mathrm i\partial\bar\partial.
\]
Weighted Hölder norms $C^{j,\alpha}_\gamma$ use the derivative
weights $\varrho_p^{l-\gamma}$ and the Hölder seminorm on
balls of radius $c\varrho_p$, with weight
$\varrho_p^{j+\alpha-\gamma}$. All charts are rescaled by
their radius; the fixed exterior uses ordinary Hölder norms.

\subsection{The domain and its normalization}

The explicit construction of the initial metric gives
\begin{equation}\label{eq:initial-local-derivatives}
 \|\widetilde g_p^R-\sigma_{i,p}\|_{C^{k+2,\alpha}_\mu
                            (\{r_i<5r_0\},\sigma_{i,p})}
       \le C.
\end{equation}
Here a tensor of weight $\mu$ has pointwise size
$O(\varrho_p^\mu)$ in the indicated metric.
This is a differentiated comparison in the actual
holomorphic node chart. It follows directly on the
core, where the difference is zero. On the radius-$R$
transition, differentiation of
\eqref{eq:initial-weight-gluing} gives the bound from
the order $2+\mu$ central potential. On
$2R<r_i<5r_0$, use the central conical expansion and
smooth convergence at fixed $R$. The latter errors
can be made at most one in this weighted norm by
shrinking the parameter neighbourhood. Higher
derivative bounds give its Hölder part. The constant
in \eqref{eq:initial-local-derivatives} is independent
of $R,p$; the permitted neighbourhood depends on $R$.

Choose a nonnegative smooth density supported in a
fixed compact subset of the unit cap and positive
on an open set. The Gram matrix of the fifteen
functions $H_a$ against this density is nonsingular:
a vanishing combination on that open set vanishes
everywhere by unique continuation, and its independent
leading terms then have zero coefficients. Let
$\ell_a$ be its dual moment functionals, so that
$\ell_a(H_b)=\delta_{ab}$.

On the $i$-th scaled cap define
\begin{equation}\label{eq:physical-cap-moments}
 H_{i,a,p}=\varepsilon_i^{d_a}H_a,\qquad
 b_{i,a}(u)=\varepsilon_i^{-d_a}\ell_a(u|_i).
\end{equation}
The right side is read after pullback to the unit cap.
Fix cutoffs $\chi_i=1$ for $r_i<r_0$ and $\chi_i=0$
for $r_i>2r_0$. Their supports are disjoint, and they
are one on the rescaled moment supports. Put
\begin{equation}\label{eq:global-domain-norm}
 v(u)=u-\sum_{i,a}b_{i,a}(u)\chi_iH_{i,a,p},
 \qquad
 \|u\|_{X_p^R}
   =\|v(u)\|_{C^{k+2,\alpha}_{2+\eta}}
                        +\sum_{i,a}|b_{i,a}(u)|.
\end{equation}
Every core moment of $v(u)$ is zero. Thus
\eqref{eq:global-domain-norm} is a unique decomposition,
equivalently a product of the closed zero-moment
remainder space and the 120 real coefficients.
For each fixed smooth fibre it is an equivalent
norm on the ordinary Hölder space; uniform
equivalence with the unweighted norm is not claimed.

Choose a fixed exterior probability density whose
support misses the cutoffs and transport it by the
exterior identifications. Let $Q(u)$ be its integral
against $u$. Define
\[
 \begin{split}
 X_{p,0}^R&=\{u\in X_p^R:Q(u)=0\},\\
 Y_p^R&=C^{k,\alpha}_\eta,\qquad
 Y_{p,0}^R=\left\{f\in Y_p^R:
                \int_{X_p}f\,\dd V_{\widetilde g_p^R}=0\right\}.
 \end{split}
\]
The domain normalization and range compatibility
use different densities.

\begin{theorem}\label{thm:global-scalar-inverse}
There are $R_1>0$ and $C$ such that, for each
$0<R<R_1$ and all sufficiently small $p\in B^\circ$,
\begin{equation}\label{eq:global-scalar-inverse}
 P_p^R:X_{p,0}^R\longrightarrow Y_{p,0}^R
 \quad\hbox{is an isomorphism},\qquad
 \|(P_p^R)^{-1}f\|_{X_p^R}\le C\|f\|_{Y_p^R}.
\end{equation}
Its forward norm is also uniformly bounded. Moreover,
\begin{equation}\label{eq:domain-hessian-control}
 \|u\|_\infty+
 \max_{j\le k}\sup\varrho_p^j
   |\nabla^j(\mathrm i\partial\bar\partial u)|
       \le C\|u\|_{X_p^R}.
\end{equation}
The constants depend on the fixed $r_0,\eta,k,\alpha$
and central geometry, but not on $R$ or the relative
cap scales.
\end{theorem}

\subsection{A negative-weight inverse and a global supremum bound}

Fix $-4<\beta<0$ and $q>6$. Pacini's
desingularization theorem
\cite[Definition~11.6 and Theorem~12.3(1)]{pacini}
applies to the original fixed central conical metric
and the eight separate complete Stenzel caps.
Its models $m_\varepsilon$ match at
$r_i\asymp\varepsilon_i^\theta$ for one fixed
$0<\theta<1$. Independent scales are allowed because
the caps are separate connected components before
gluing. The central space is connected, so the
kernel consists of one global constant. Its
complement can be chosen as $\ker Q$: the extension
cutoff equals one on the fixed support of $Q$.
The resulting negative-weight inverse constant is
fixed before $R$ is chosen.

We compare these models with the initial metric
through scaled derivatives:
\begin{equation}\label{eq:pacini-model-comparison}
 \max_{j\le k+2}\sup\varrho_p^j
 |\nabla_{m_\varepsilon}^j
       (\Phi_p^*\widetilde g_p^R-m_\varepsilon)|
 \le CR^\mu+e_R(p),\qquad e_R(p)\longrightarrow0.
\end{equation}
For fixed $R$, the gluing radii lie inside $R$.
There the initial metric is Stenzel; at the model
transition the conically singular and asymptotically
conical errors tend to zero separately with each
$\varepsilon_i$. Between that transition and $R$,
the difference from the central metric is bounded by
$CR^\mu$ and the cap tail. The initial transition
and the exterior obey the same estimate.
For phases choose $\theta<a<1$ and interpolate any
bounded Reeb rotation between
$\varepsilon_i^a$ and $\varepsilon_i^\theta$.
Every scaled derivative of a smooth logarithmic
interpolation is $O(1/|\log\varepsilon_i|)$.
Both endpoints are asymptotically conical, since
$\varepsilon_i^a/\varepsilon_i\to\infty$.
This proves \eqref{eq:pacini-model-comparison},
including phases and differentiated operator
coefficients, without a ratio assumption.

For the possibly non-Kähler model let $P_m$ be
the same fixed multiple of the real scalar
Laplacian as $P_p^R$ on the Kähler fibre.
The augmented model map
\[
 (u,c)\longmapsto P_m u+c,\qquad Q(u)=0,
\]
is a uniform isomorphism from
$(W^{k+2,q}_\beta\cap\ker Q)\oplus\RR$
to $W^{k,q}_{\beta-2}$. Indeed $c$ is the
model-volume mean of the target. Its norm is
uniformly bounded by Hölder's inequality in the
measure $\varrho^{-6}\dd V$: the required radial
power is $\varrho^{q'(\beta+4)}$, which is
integrable since $\beta>-4$. The total volume
has a positive exterior lower bound, and constant
functions have uniformly bounded domain and target
norms at these negative weights. Pacini's
normalized inverse handles the remaining mean-zero
target.

By \eqref{eq:pacini-model-comparison}, the two
augmented maps differ by a small operator norm,
first for small fixed $R$ and then small $p$.
Neumann inversion gives the corresponding bound
for the initial metric. If $f$ has zero current
volume mean, integrating $P_p^Ru+c=f$ gives $c=0$.
Consequently
\begin{equation}\label{eq:negative-weight-inverse}
 \|u\|_{W^{k+2,q}_\beta}
 \le C\|f\|_{W^{k,q}_{\beta-2}}
 \le C\|f\|_{Y_p^R},\qquad Q(u)=0.
\end{equation}
The second inequality uses the positive radial
power $\eta-\beta+2$ and has no logarithmic loss.
Weighted embedding bounds $u$ on the complement
of the radius-$r_0$ regions.

To control the cores, use the exact Stenzel
potential $\phi_{i,p}=\varepsilon_i^2\phi$ as
a barrier. It is bounded by $Cr_0^2$ there, and
\eqref{eq:initial-local-derivatives}, with $r_0$
small, gives
\[
 P_p^R\phi_{i,p}
       =\tr_{\widetilde\omega_p^R}\sigma_{i,p}\ge c_*>0.
\]
Apply the maximum principle to
$u\pm c_*^{-1}\|f\|_\infty\phi_{i,p}$ with
boundary at radius $r_0$. This gives
\begin{equation}\label{eq:uniform-potential-supremum}
 \|u\|_\infty\le C(r_0)\|f\|_{Y_p^R}.
\end{equation}
Ordinary interior Schauder estimates now control
all required norms on fixed-radius annuli.

\subsection{The positive-weight cap expansion}

On the fixed complete Stenzel cap,
\cite[Theorems~2.9 and 2.11]{conlon-hein}
gives a surjective Fredholm map
$C^{k+2,\alpha}_{2+\eta}\to C^{k,\alpha}_\eta$.
Its kernel is Lemma~\ref{lem:cap-harmonic-basis}.
Subtract the core moments from a bounded right
inverse to obtain $G_\delta$ with
$\ell_a(G_\delta F)=0$.

For completeness, the Hölder seminorm printed
in that reference compares pairs below a fixed
convexity radius. Our scaled input norm controls
that norm, so the inverse first gives the weighted
supremum estimate for its output. Schauder
estimates on balls comparable to $\rho$ then give
the full scaled Hölder estimate, with the compact
core treated by fixed balls. Thus this inverse
is bounded in precisely our norms, with no loss
in $\eta$. Under homothety its physical norm
remains uniform because $\delta=\eta+2$.

Choose $\zeta=1$ for $r_i<3r_0$ and $\zeta=0$
for $r_i>4r_0$. On the complete scaled cap set
\begin{equation}\label{eq:global-local-expansion}
 w=G_{\delta,i,p}P_{\sigma_{i,p}}(\zeta u),
 \qquad
 H=\zeta u-w=\sum_ab_aH_{i,a,p}.
\end{equation}
Since $w$ has zero core moments, these $b_a$
are exactly \eqref{eq:physical-cap-moments}.
Put $V_i=\|w\|_{C^{k+2,\alpha}_\delta}$ and
$U=\|u\|_\infty$. The annular Gram estimate
\eqref{eq:annular-harmonic-coefficients}, now at
physical radius $r_0$, gives
\begin{equation}\label{eq:global-quadratic-coefficients}
 |b_{(2)}|\le C(r_0^{-2}U+r_0^\eta V_i)
\end{equation}
for the six degree-two coefficients.

On $r_i<3r_0$, the source in
\eqref{eq:global-local-expansion} is
$f+(P_\sigma-P_p^R)u$, and $u=H+w$.
The constant and coordinate components of $H$
have identically zero complex Hessian for both
Kähler metrics. For the other components,
\eqref{eq:initial-local-derivatives} yields
\[
 \begin{split}
 \|(P_\sigma-P_p^R)H\|_{C^{k,\alpha}_\eta}
       &\le Cr_0^{\mu-\eta}|b_{(2)}|,\\
 \|(P_\sigma-P_p^R)w\|_{C^{k,\alpha}_\eta}
       &\le Cr_0^\mu V_i.
 \end{split}
\]
The use of the complex Hessian is essential here;
the real Hessians of the pluriharmonic coordinate
functions do not vanish. On the cutoff annulus,
\eqref{eq:uniform-potential-supremum} and ordinary
Schauder estimates instead bound all source and
commutator terms. The cap inverse therefore gives
\[
 V_i\le C\left(\|f\|_Y+r_0^{-\delta}U
             +r_0^\mu V_i+r_0^{\mu-\eta}|b_{(2)}|\right).
\]
Insert \eqref{eq:global-quadratic-coefficients}
and fix $r_0$ small enough to absorb the resulting
$Cr_0^\mu V_i$. Equation
\eqref{eq:uniform-potential-supremum} gives
$V_i\le C(r_0)\|f\|_Y$. The same annular Gram
estimate then bounds all fifteen coefficients,
including degrees zero and $3/2$. On $r_i<r_0$
the canonical remainder $v(u)$ equals $w$.
On the fixed cutoff annuli and exterior its
norm is controlled by the coefficient bounds
and ordinary Schauder estimates. This proves
the inverse estimate in
\eqref{eq:global-scalar-inverse}.

\subsection{Isomorphism and the tensor estimate}

For each fixed smooth compact fibre, the scalar
equation has a solution for every zero-volume-mean
Hölder target, unique up to a constant. Impose
$Q(u)=0$. The preceding estimates apply directly
to these Hölder solutions by elliptic theory;
no assertion of density of smooth functions in
the full Hölder norm is needed.

The forward map is also bounded. The remainder
has weight $2+\eta$, so its image has weight
$\eta$. On a node chart the pluriharmonic modes
have zero image, and the quadratic modes have
image $O(\varrho_p^\mu)$ by
\eqref{eq:initial-local-derivatives}. Cutoff
derivatives occur at fixed radius. Integration
gives the required range condition, and the
only global harmonic functions are constants,
removed by $Q$.

Finally the remainder has complex Hessian
$O(\varrho_p^\eta)$; each quadratic mode has
bounded complex Hessian with scaled derivatives;
and each pluriharmonic mode has zero such
Hessian where its cutoff is constant. All other
terms occur at fixed radius. Their supremum
norms are bounded as well. This proves
\eqref{eq:domain-hessian-control} and completes
the proof of Theorem~\ref{thm:global-scalar-inverse}.
